\subsection{What the growth and zeros count}
\label{grw:section}

The operator behind $\Psi_g$ controls the location of its zeros.
The companion $\Phi_g$ records how quickly disabled lanes accumulate.
If those primes grow sparse, the function grows slowly; if they grow
more numerous, it grows faster.  For the Goldbach member this comparison
also turns the conjecture into a statement about complex zeros.

Put $a_k=M_k(g)/k!$ and $B_g(R)=\#(\mathcal F\cap[2,R])$.
If all $k$ labels in the prime-assignment model of
Proposition~\ref{fnd:moment-positivity} are masked, their first
assigned primes are $k$ distinct members of $\mathcal F$.
For a chosen set of $k$ primes, sum over its $k!$ assignments to the
labels and discard the first-assignment restrictions.  Independence
between primes gives the useful bound
\begin{equation}
 a_k\le e_k\left(\frac1{p+k-1}:p\in\mathcal F\right)
 \le e_k\left(\frac1p:p\in\mathcal F\right),
 \label{grw:coefficient-majorant}
\end{equation}
Here $e_k$ means the sum of products over all $k$-element subsets; for
example $e_2(w)=\sum_{p<q}w_pw_q$.  Thus, when
$\sum_{p\in\mathcal F}1/p<\infty$,
\begin{equation}
 |\Phi_g(z)|\le\Phi_g(|z|)
 \le P_{\mathcal F}(|z|):=
       \prod_{p\in\mathcal F}(1+|z|/p).
 \label{grw:majorant}
\end{equation}
The product bounds the function from above.  Its factors are not being
identified with the zeros of $\Phi_g$.

The \emph{order} of an entire function records the power in its
exponential growth: roughly, growth like $\exp(r^\alpha)$ on disks of
radius $r$ gives order $\alpha$.  The following theorem matches that
analytic exponent with a counting exponent for the disabled primes.

\begin{theorem}[Counting exponent and entire-function order]
\label{grw:order}
Suppose the disabled-prime counting exponent
\[
 \eta_g=\limsup_{R\to\infty}
 \frac{\log\max(1,B_g(R))}{\log R}
\]
is strictly less than one.  Then $\Phi_g$ has order exactly $\eta_g$.
Formally the order is the limsup of $\log\log \max_{|z|\le r}|\Phi_g(z)|$
divided by $\log r$; constants and polynomials have order zero by convention.  For the actual
Goldbach mask,
\begin{equation}
 \operatorname{ord}\Phi_{\Gold}
 =\limsup_{Y\to\infty}\frac{\log\max(1,E(Y))}{\log Y}
 \le\frac{18}{25},
 \label{grw:goldbach-order}
\end{equation}
and, as $r\to\infty$,
\begin{equation}
 \log\max_{|z|\le r}|\Phi_{\Gold}(z)|
 =\log\Phi_{\Gold}(r)\ll(r/\log r)^{18/25}.
 \label{grw:quantitative-growth}
\end{equation}
In particular, for an absolute constant $C$ and $k\ge2$,
\begin{equation}
 M_k(\Gold)\le
 \left(\frac{C}{k^{7/18}\log(k+1)}\right)^k.
 \label{grw:moment-bound}
\end{equation}
\end{theorem}
The proof is in Appendix~\ref{grw:proofs}.  Its two estimates have
opposite roles.  Summing over every possible set of disabled primes
bounds the function from above.  Requiring the first $k$ disabled
primes to hit distinct labels supplies a matching lower bound on its
order.  The analytic exponent therefore cannot lose the arithmetic
counting exponent.


\begin{corollary}[Zeros count the exceptions]\label{grw:zeros}
Under the hypothesis of Theorem~\ref{grw:order}, list the zeros
of $\Phi_g$ with multiplicity as $(\zeta_j)$.  Then
\begin{equation}
 \Phi_g(z)=\prod_j(1-z/\zeta_j),\qquad
 \sum_j|\zeta_j|^{-1}<\infty,\qquad
 \sum_j\zeta_j^{-1}=-S_g.
 \label{grw:factorization}
\end{equation}
Counting multiplicity, the number of zeros equals the number of
disabled lanes, allowing infinity.
If $n_g(R)$ counts the zeros in $|z|\le R$, with multiplicity, then
\begin{equation}
 \limsup_{R\to\infty}
 \frac{\log\max(1,n_g(R))}{\log R}=\eta_g.
 \label{grw:zero-exponent}
\end{equation}
For the actual Goldbach function,
\[
 \GC\ \Longleftrightarrow\ \Phi_{\Gold}\text{ has no complex zeros},
 \qquad n_{\Gold}(R)\ll(R/\log R)^{18/25}.
\]
If there are infinitely many exceptions, every complex value is
attained by $\Phi_{\Gold}$ infinitely often.
\end{corollary}
\begin{proof}
Hadamard's factorization theorem \cite{sampHadamard1893} says that
an entire function of order less than one has a genus-zero product
and a constant exponential prefactor.  Since $\Phi_g(0)=1$, this
gives the first two assertions of \eqref{grw:factorization};
differentiation at zero gives its last assertion because
$\Phi_g'(0)=M_1(g)=S_g$.  For a finite mask the polynomial degree
is $\#\mathcal F$.  For an infinite mask all coefficients are
positive, so the function is not a polynomial; the factorization
then forces infinitely many zeros.  The same argument applied to
$\Phi_g-w$ shows infinite attainment of each $w\in\mathbb C$
when the mask is infinite.

Jensen's formula and $\Phi_g(0)=1$ give
$n_g(R)\log2\le\log\Phi_g(2R)$, which bounds the zero-counting
exponent by $\eta_g$ and gives the displayed quantitative estimate.
Conversely, if that exponent were $\lambda<\eta_g$, choose
$\lambda<\beta<\eta_g<1$.  The product
\eqref{grw:factorization}, together with
$n_g(t)\ll_\beta t^\beta$, implies
\[
 \log\max_{|z|\le r}|\Phi_g(z)|
 \le\sum_j\log(1+r/|\zeta_j|)=O_\beta(r^\beta)
\]
by the same integration used in \eqref{grw:stieltjes}, starting
below the smallest zero modulus.  This contradicts its order.
The zero-order case follows from the upper bound and nonnegativity.
\end{proof}

These are zeros of the mask function.  The worked example above shows
that they can be nonreal.  The sparsity hypothesis is essential to the
factorization argument.
Also, order zero alone does not mean finitely many exceptions: an
infinite mask with subpower counting function has the same order.
The improved moment estimate measures a statistic of our
construction; it does not improve the imported bound for $E(Y)$.

\subsection{One bias can hide infinitely many changes}
\label{grw:collision-section}

There is a sharp reason to keep the higher correlations rather
than only the proportion of zero digits.  Put
$\delta_p=p^{-1}\prod_{q<p}(1-1/q)$, so
$S_g=\sum_{p\in\mathcal F}\delta_p$.

\begin{proposition}\label{grw:bias-collision}
Among masks known to have finitely many disabled lanes, the exact bias $S_g$ determines the mask.
Nevertheless, for every sufficiently large prime $p$, there is an
infinite set $\mathcal H=\{q_1<q_2<\cdots\}$ of primes greater than
$p$ such that
\begin{equation}
 \sum_{q\in\mathcal H}\delta_q=\delta_p,\qquad
 q_{j+1}/q_j\longrightarrow\infty,
 \qquad \#(\mathcal H\cap[2,R])=o(\log R).
 \label{grw:collision}
\end{equation}
Thus two masks can have exactly the same digit bias while one has
one disabled lane and the other has infinitely many.
\end{proposition}
\begin{proof}
If $P$ is the largest prime of a finite nonempty mask, then
$\delta_P$ has $P$-adic valuation $-1$: its numerator contains only
factors $q-1<P$, and its denominator is the primorial through $P$.
Every earlier $\delta_q$ has nonnegative $P$-adic valuation.
Consequently $P$ is the largest prime dividing the reduced
denominator of $S_g$.  Recover $P$, subtract $\delta_P$, and repeat.

List all primes as $p_i$ and set $d_i=\delta_{p_i}$.  These weights
decrease strictly to zero and
\[
 d_i/d_{i+1}=p_{i+1}/(p_i-1)\longrightarrow1
\]
by the prime number theorem.  Choose $p$ beyond a point where every
consecutive ratio is less than two.  Starting with residual
$r_0=\delta_p$, repeatedly subtract the largest weight at a prime
greater than $p$ which does not exceed the residual.
The first choice is the successor of $p$.  If $q$ is a later
choice and $q^-$ its immediate prime predecessor, the residual
after subtraction satisfies
\[
 0\le r_{\mathrm{new}}
 \le\delta_{q^-}-\delta_q=o(\delta_q).
\]
The ratio bound makes $r_{\mathrm{new}}<\delta_q$, so no weight is
used twice and selected primes strictly increase.  The process
cannot terminate: a finite sum at primes greater than $p$ has a
reduced denominator containing its largest prime, unlike
$\delta_p$.  It therefore makes infinitely many selections, and
$r_j<\delta_{q_j}\to0$ proves the sum identity.

The next selected weight is at most the preceding residual, hence
$\delta_{q_{j+1}}/\delta_{q_j}\to0$.  Mertens' product estimate
$\delta_q\sim e^{-\gamma}/(q\log q)$ forces
$q_{j+1}/q_j\to\infty$.  For any fixed $C>1$, these ratios
eventually exceed $C$, so the counting function is at most
$\log R/\log C+O_C(1)$.  Letting $C$ tend to infinity proves
\eqref{grw:collision}.
\end{proof}

The equality is exact, so it survives arbitrarily accurate measurement
of the digit frequency.  For these two masks $\Phi_g'(0)$ agrees,
but their functions are
respectively linear and transcendental entire.  The second has
order zero and infinitely many zeros by the preceding results.
This is an example within the mask family; it asserts no actual
Goldbach exception, and no collision between the whole functions.
